% MSR 2027 - Data and Tool Showcase Track
% 4 pages + 1 page of references. Single-anonymous: author names ARE listed.
% Submission: https://msr2027-data-tool.hotcrp.com/
% Abstract deadline 2026-11-05 AoE, paper deadline 2026-11-10 AoE.
\documentclass[10pt,conference]{IEEEtran}

\usepackage{booktabs}
\usepackage{graphicx}
\usepackage{listings}
\usepackage{url}
\usepackage{xcolor}

% Numbers still being regenerated. Every \pending must be gone before
% submission - grep for it.
\newcommand{\pending}[1]{\textcolor{red}{#1}}

\newcommand{\dsname}{CRUSTPairs}   % TODO: confirm the dataset name

\begin{document}

\title{\dsname: A Cross-Language Dataset of Paired C and Rust Solutions}

\author{
\IEEEauthorblockN{Leonardo J. Almeida}
\IEEEauthorblockA{University of S\~ao Paulo\\
leonardo.j.almeida@usp.br}
\and
\IEEEauthorblockN{Vojt\v{e}ch Hejlek}
\IEEEauthorblockA{TODO affiliation\\
TODO email}
% TODO: confirm author list, order, and the advisor's inclusion.
}

\maketitle

\begin{abstract}
TODO. Roughly: C-to-Rust migration is an active agenda, but no dataset pairs
the two languages at the level of whole solutions to the same problem. We
present \dsname, 1857 C/Rust implementation pairs drawn from competitive
programming archives, normalised, deduplicated, and stratified into difficulty
tiers by zero-shot embedding similarity. We report retrieval baselines showing
that a generic code model finds these pairs hard zero-shot and that the gap
closes substantially with finetuning.
\end{abstract}

\begin{IEEEkeywords}
dataset, cross-language, code similarity, Rust, C, code retrieval
\end{IEEEkeywords}

\section{Introduction}
\label{sec:intro}

% ~0.5 page. Why C<->Rust specifically:
%  - memory-safety driven migration is an active agenda, including public-sector
%    pushes; C-to-Rust translation tools (c2rust, LLM-based) are proliferating
%  - evaluating them needs paired ground truth, which does not exist at the
%    level of whole solutions to a shared problem
%  - existing cross-language corpora stop short of Rust (see Sec. II)
% Close with the contribution list.

\noindent\textbf{Contributions.}
\begin{itemize}
  \item \dsname, 1857 C/Rust solution pairs with problem statements,
        normalised entry points, per-record provenance and licensing.
  \item A difficulty stratification derived from zero-shot cross-language
        embedding similarity, independent of the model family used in
        evaluation.
  \item Retrieval baselines, per difficulty tier, over a frozen stratified
        split.
\end{itemize}

\section{Related Datasets}
\label{sec:related}

% ~0.5 page. The positioning table is the argument: a "covers Rust?" column
% with a column of No.
% Cover: CodeNet, xCodeEval, XLCoST, CoST, AVATAR, CodeXGLUE/POJ-104.
% Say explicitly what each does and where it stops.

\begin{table}[t]
\caption{Cross-language code corpora. \dsname{} is the only one pairing C with
Rust at solution granularity. TODO: verify every cell.}
\label{tab:related}
\centering
\begin{tabular}{lccc}
\toprule
Dataset & Languages & Paired & Rust \\
\midrule
CodeNet        & 55  & no  & no \\
xCodeEval      & 17  & yes & no \\
XLCoST         & 7   & yes & no \\
CoST           & 7   & yes & no \\
AVATAR         & 2   & yes & no \\
\midrule
\dsname{}      & 2   & yes & yes \\
\bottomrule
\end{tabular}
\end{table}

\section{Construction}
\label{sec:construction}

% ~0.75 page. Figure 1 is the pipeline with counts at each stage.

\subsection{Sources and pairing}

\dsname{} draws on two competitive programming archives: xCodeEval
\cite{khan2023xcodeeval}, itself derived from Codeforces, and IBM Project
CodeNet \cite{puri2021codenet}, derived from AtCoder. Both distribute accepted
submissions, so the two implementations of a pair are known to have passed the
same official tests. A pair is one problem, one accepted C submission, and one
accepted Rust submission.

\subsection{Normalisation}

Both sides pass through a static analysis pipeline that strips comments,
removes unused functions and macros, normalises encoding, and applies the
language's standard formatter (\texttt{clang-format}, \texttt{rustfmt}). Entry
points are renamed to \texttt{solution} in both languages so that the function
name carries no signal. One consequence, discussed in
Section~\ref{sec:limitations}, is that the distributed code is no longer a
standalone compilable program.

\subsection{Cleaning}

From an initial 2013 pairs, 127 were removed: 2 confirmed mis-pairs, 69 AtCoder
easy/hard variants whose statements were near-identical (description embedding
similarity $>0.90$), 54 further suspicious pairs ($0.80$--$0.90$), 1 length
ratio outlier, and 1 case of genuine implementation divergence. Near-duplicate
detection used OpenAI \texttt{text-embedding-3-large} over the raw statements.

A further 29 pairs, a set of textbook reference implementations, were dropped
for two converging reasons. They carried the dataset's only copyleft content
and were its only records whose two halves were under different licences
(Section~\ref{sec:availability}); and they were structurally unlike the rest,
having no problem statement and functions taking arguments rather than reading
standard input. This leaves 1857 pairs.

\subsection{Difficulty stratification}

Each pair is scored by the cosine similarity between zero-shot
\texttt{SFR-Embedding-Code-400M\_R} embeddings of its C and Rust sides, and
binned at the 25th and 75th percentiles into easy, medium and hard. Pairs a
generic code model already embeds close together are near-literal translations;
pairs it separates diverge structurally and demand a real cross-language
remapping.

The choice of SFR rather than UniXcoder is deliberate. Defining the tiers with
the same model family that is later finetuned would make hard pairs hard by
construction, and any reported per-tier improvement would be circular. SFR is
trained on a different corpus, so the tiers remain an independent axis. The
stratification is also origin-agnostic: AtCoder and Codeforces publish
difficulty ratings on incompatible scales, so upstream ratings are not used.

\section{Dataset Description}
\label{sec:description}

\subsection{Schema}

% TODO: schema table. Fields: problem_id, origin, license,
% problem_description, problem_description_format, c_code, rust_code,
% difficulty, difficulty_score, split.

\subsection{Statistics}

\begin{table}[t]
\caption{Size of the two sides, in characters and lines.}
\label{tab:sizes}
\centering
\begin{tabular}{lrrrr}
\toprule
 & Mean & Median & p90 & Max \\
\midrule
C, characters     & 616  & 516  & 1046 & 16460 \\
Rust, characters  & 851  & 749  & 1424 & 4888 \\
Statement, chars  & 1459 & 1231 & 2717 & 7368 \\
\midrule
C, lines          & 33.6 & & & \\
Rust, lines       & 31.4 & & & \\
\bottomrule
\end{tabular}
\end{table}

The Rust side is about $1.5\times$ longer than the C side in characters
(median ratio $1.50$) while being \emph{shorter} in lines (31.4 against 33.6):
Rust solutions pack more into each line. This is worth stating because it means
a length-based heuristic does not separate the two languages the way a naive
reading of the character counts would suggest.

\begin{table}[t]
\caption{Composition by source and difficulty tier.
\pending{Tier counts pending regeneration over the 1857-pair population.}}
\label{tab:composition}
\centering
\begin{tabular}{lrrrr}
\toprule
Source & Easy & Medium & Hard & Total \\
\midrule
CodeNet (AtCoder)     & \pending{251} & \pending{403} & \pending{185} & 839 \\
xCodeEval (Codeforces)& \pending{207} & \pending{528} & \pending{283} & 1018 \\
\midrule
Total                 & \pending{458} & \pending{931} & \pending{468} & 1857 \\
\bottomrule
\end{tabular}
\end{table}

Source and difficulty are not independent: the CodeNet portion skews easy while
the xCodeEval portion skews hard. Reporting a per-tier result therefore
requires a split that controls for both.

\subsection{Splits}

The train/validation/test assignment is frozen in the data rather than
regenerated from a seed, and is stratified by source $\times$ difficulty at
80/10/10: 1485, 185 and 187 pairs. The test set mirrors the full dataset to
within 0.3 percentage points by source and 0.6 by tier.

\section{Baselines}
\label{sec:baselines}

% ~0.5 page, one table. Zero-shot UniXcoder vs LoRA vs full finetune,
% MRR@10 / R@1 / R@5, overall and per tier.
% Frame as evidence of utility, NOT as the contribution.

\begin{table}[t]
\caption{Cross-language retrieval on the 187-pair test set.
\pending{All numbers pending re-run against the current split.}}
\label{tab:baselines}
\centering
\begin{tabular}{lccc}
\toprule
Model & MRR@10 & R@1 & R@5 \\
\midrule
UniXcoder, zero-shot   & \pending{--} & \pending{--} & \pending{--} \\
UniXcoder + LoRA       & \pending{--} & \pending{--} & \pending{--} \\
UniXcoder, full        & \pending{--} & \pending{--} & \pending{--} \\
\bottomrule
\end{tabular}
\end{table}

\section{Use Cases}
\label{sec:usecases}

% ~0.25 page:
%  - evaluating C-to-Rust translators (c2rust, LLM-based) against real pairs
%  - cross-language code search and clone detection (type-4, across languages)
%  - training cross-language retrieval models
%  - studying how idiom differs between a systems language with and without
%    ownership

\section{Limitations}
\label{sec:limitations}

\textbf{Equivalence is inherited, not verified.} No test cases or expected I/O
are distributed. That both submissions were accepted upstream is evidence of
functional equivalence, but this dataset does not itself execute anything, and
the entry-point renaming means the distributed code does not compile as a
standalone program. Execution-based verification is the clearest line of future
work.

\textbf{Upstream identifiers are absent.} The pairs cannot be traced back to
the originating problem or submission, which limits attribution to the corpus
level and prevents joining against upstream metadata.

\textbf{Competitive programming is not production code.} Solutions are short,
single-file, and I/O-driven. Conclusions drawn here should not be assumed to
transfer to systems code, which is where C-to-Rust migration actually matters.

\textbf{Difficulty is a proxy.} The tiers measure embedding similarity, not
algorithmic difficulty or translation effort.

\textbf{Pre-training contamination is possible.} Both source archives are
public and may appear in the pre-training corpora of the models evaluated here.

\section{Availability and Licensing}
\label{sec:availability}

% TODO: Zenodo DOI (required at camera-ready; get it before submitting),
% repository URL, Hugging Face URL.

\dsname{} aggregates two corpora under different licences, so it is not
distributed under a single one; every record carries its own SPDX expression.
The 1018 xCodeEval-derived pairs are CC BY-NC 4.0 and the 839 CodeNet-derived
pairs are CDLA-Permissive-2.0. The aggregate is therefore non-commercial, and
attribution to both upstream projects is required. Filtering to the CodeNet
portion yields a 839-pair subset with no commercial restriction.

\section*{Acknowledgements}
TODO.

\bibliographystyle{IEEEtran}
\bibliography{refs}

\end{document}
